\section{PC-RCO观测器设计}

PC-RCO的核心思想是利用CFD预标定动力学模型进行一步超前速度预测，通过灵敏度分析将预测残差映射为海流估计梯度，并以连续GM正则化和自适应限幅保证估计的鲁棒性和物理可接受性。

\subsection{速度预测架构}

观测器在第$k$步使用当前海流估计$\hat{\bm{c}}_k$对$k+1$步的AUV速度进行一步向前预测：

\begin{equation}
    \hat{\bm{\nu}}_{k+1|k} = \bm{\nu}_k + \Delta t \cdot \bm{a}_{\text{model}}(\bm{\nu}_k, \bm{\tau}_k, \hat{\bm{c}}_k)
    \label{eq:prediction}
\end{equation}

其中$\bm{a}_{\text{model}}$为CFD动力学模型计算的加速度。与加速度残差方案不同，速度预测利用时间积分的平滑效应，避免了对$\dot{\bm{\nu}}$的直接测量或差分——后者在典型DVL噪声（0.01 m/s）和100 Hz采样率下的信噪比仅约0.03，难以用于梯度估计。

\subsection{灵敏度分析}

定义速度预测对海流参数的灵敏度矩阵：

\begin{equation}
    \bm{\Phi}_c = \frac{\partial \hat{\bm{\nu}}_{k+1|k}}{\partial \bm{c}} \in \R^{6\times 2}
    \label{eq:sensitivity}
\end{equation}

由于CFD阻力模型$\bm{D}(\bm{\nu}_r)$具有解析的二次形式，其对$\bm{c}$的偏导数存在闭合形式。然而，完整6自由度动力学方程的解析求导较为繁琐，因此本文采用数值扰动法计算$\bm{\Phi}_c$：

\begin{equation}
    \bm{\Phi}_c(:,j) = \frac{\hat{\bm{\nu}}_{k+1|k}(\hat{\bm{c}}+\delta\bm{e}_j) - \hat{\bm{\nu}}_{k+1|k}(\hat{\bm{c}})}{\delta}, \quad j=1,2
    \label{eq:numerical_phi}
\end{equation}

其中扰动步长$\delta=0.01$ m/s，$\bm{e}_j$为单位向量。每次更新需要3次CFD模型前向计算（1次基准+2次扰动），计算开销可控。

灵敏度Gramian定义为：

\begin{equation}
    \bm{W}_c = \bm{\Phi}_c^T\bm{\Phi}_c \in \R^{2\times 2}
    \label{eq:gramian}
\end{equation}

$\bm{W}_c$近似于海流参数估计的Fisher信息矩阵，其最小特征值$\lambda_{\min}(\bm{W}_c)$定量刻画当前动力学状态下海流参数的可辨识程度。

\subsection{连续Gauss-Markov正则化}

PC-RCO摒弃了二值激励门控（$\lambda_{\min}(\bm{W}_c) > \mu_{\text{gate}}$时更新、否则GM传播），而是在每个时间步同时施加梯度更新和GM衰减：

\begin{equation}
    \hat{\bm{c}}_{k+1} = e^{-\Delta t/\tau_c}\left(\hat{\bm{c}}_k + \Delta\bm{c}_k^{\text{grad}}\right) + (1-e^{-\Delta t/\tau_c})\bar{\bm{c}}
    \label{eq:continuous_gm}
\end{equation}

其中$\Delta\bm{c}_k^{\text{grad}}$为梯度驱动更新（第3.4节），$\bar{\bm{c}}$为先验均值。该连续公式具有两种工作模式：
\begin{itemize}[itemsep=0pt,topsep=4pt]
    \item \textbf{高激励}：$\|\Delta\bm{c}_k^{\text{grad}}\| \gg (1-e^{-\Delta t/\tau_c})\|\hat{\bm{c}}_k - \bar{\bm{c}}\|$——梯度更新主导，GM衰减作为弱Tikhonov正则化。
    \item \textbf{低激励}：$\|\Delta\bm{c}_k^{\text{grad}}\| \ll (1-e^{-\Delta t/\tau_c})\|\hat{\bm{c}}_k - \bar{\bm{c}}\|$——GM衰减主导，逐渐将估计拉向先验，防止漂移。
\end{itemize}

\textbf{连续优于二值门控的动机}：通过系统诊断，我们发现基于灵敏度Gramian特征值的二值激励门控在反馈控制AUV中永远保持开放状态。根本原因有二：（1）CFD阻力模型的二次型在极低转弯速率下仍提供非零灵敏度；（2）反馈控制器（如ALOS制导）持续产生小幅度航向修正以对抗横向海流漂移，阻止持续零激励条件。连续方案消除了阈值调谐需求，在任何激励水平下提供鲁棒正则化。

\subsection{自适应限幅的梯度更新}

式\eqref{eq:continuous_gm}中的梯度驱动更新$\Delta\bm{c}_k^{\text{grad}}$计算如下：

\begin{equation}
    \Delta\bm{c}_k^{\text{grad}} = \gamma \cdot \frac{\bm{\Phi}_c^T \bm{e}_{\nu}}{1 + \lambda_{\min}(\bm{W}_c) \cdot \kappa}
    \label{eq:update}
\end{equation}

其中$\bm{e}_{\nu} = \bm{\nu}_{k+1}^{\text{meas}} - \hat{\bm{\nu}}_{k+1|k}$为速度预测新息（仅取纵荡和横荡分量），$\gamma = K_{\text{obs}} \cdot 100$为有效增益，$\kappa$为自适应正则化系数。选择梯度方向$\bm{\Phi}_c^T\bm{e}_{\nu}$而非Newton方向，是为了避免Gramian近奇异时的数值不稳定，并在模型失配下提供更强的鲁棒性。

\textbf{基于在线噪声估计的自适应限幅}：为平衡收敛速度与物理可接受性，采用动态调整的单步限幅：

\begin{equation}
    \|\Delta\bm{c}_k^{\text{grad}}\| \leq \Delta c_{\max}(\hat{\sigma}_k), \quad \Delta c_{\max}(\hat{\sigma}_k) = \min\left(\Delta c_{\max}^{\text{base}} \cdot \max\left(0.5,\; \min\left(3.0,\; \frac{\hat{\sigma}_k}{\sigma_{\text{ref}}}\right)\right),\; 0.20\right)
    \label{eq:adaptive_clamp}
\end{equation}

其中$\hat{\sigma}_k$为从速度预测残差计算的EMA噪声估计：

\begin{equation}
    \hat{\sigma}_k = 0.995\,\hat{\sigma}_{k-1} + 0.005\,\|\bm{e}_{\nu}\|
    \label{eq:noise_est}
\end{equation}

$\sigma_{\text{ref}}=0.02$ m/s为参考DVL噪声水平，$\Delta c_{\max}^{\text{base}}=0.06$ m/s为名义限幅值。自适应方案：（1）在正常低噪声条件下维持紧限幅（$\approx$0.06 m/s）保证平滑收敛；（2）在高噪声下放宽至约0.18 m/s以维持对噪声的穿透能力；（3）强制执行0.20 m/s的绝对上限，防止任何条件下的无界更新。

\subsection{前馈补偿策略}

海流估计$\hat{\bm{c}}$最终用于前馈补偿。将$\hat{\bm{c}}$转化为体坐标系海流分量，计算海流引起的净力/力矩效应：

\begin{equation}
    \hat{\bm{d}} = \bm{M}\left[\dot{\bm{\nu}}(\hat{\bm{c}}) - \dot{\bm{\nu}}(\bm{0})\right]
    \label{eq:feedforward}
\end{equation}

在XHY AUV的5推进器配置中，推力分配矩阵的第4行（横滚通道）全为0，5推进器无法产生独立横滚力矩。补偿消融实验发现，纵荡前馈补偿通过推力分配矩阵间接影响偏航力矩，产生纵荡-偏航耦合。基于此，PC-RCO采用仅偏航前馈策略：仅将$\hat{\bm{d}}(6)$（偏航力矩补偿）注入控制律。
